
\exo

Sans calculatrice, donner la valeur des nombres suivants.

\vspace{-6pt}
\setlength{\columnsep}{1cm}
\setlength{\columnseprule}{0pt}
\begin{multicols}{3}
\begin{enumerate}
\item $a=\big(\sqrt{25}\big)^2$.
\medskip
\item $b=\sqrt{3^2}$.

\item $c=\big(-\sqrt{16}\big)^2$.
\medskip
\item $d=\big(\sqrt{0,14}\big)^2$.

\item $e=\sqrt{(-7)^2}$.
\medskip
\item $f=\big(\sqrt{0,4^2}\big)^2$.
\end{enumerate}
\end{multicols}
\vspace{-6pt}

\exo

Écrire les nombres suivants sous la forme $a\sqrt{3}$, où $a$ est un entier.

\vspace{-6pt}
\setlength{\columnsep}{1cm}
\setlength{\columnseprule}{0pt}
\begin{multicols}{3}
\begin{enumerate}
\item $a=\sqrt{5}\times\sqrt{15}$.
\item $b=\sqrt{75}$.
\item $c=\sqrt{108}$.
\end{enumerate}
\end{multicols}
\vspace{-6pt}

\exo

Écrire les nombres suivants sous la forme $a\sqrt{b}$, où $a$ et $b$ sont des entiers, $b$ étant le plus petit possible.

\vspace{-6pt}
\setlength{\columnsep}{1cm}
\setlength{\columnseprule}{0pt}
\begin{multicols}{3}
\begin{enumerate}
\item $a=\sqrt{32}$.
\medskip
\item $b=\sqrt{75}$.
\medskip
\item $c=\sqrt{500}$.

\item $d=\sqrt{80}$.
\medskip
\item $e=5\sqrt{18}$.
\medskip
\item $f=4\sqrt{32}$.

\item $g=\sqrt{2}\times \sqrt{6}$.
\medskip
\item $h=\sqrt{7}\times 3\sqrt{14}$.
\medskip
\item $k=7\sqrt{2}\times 5\sqrt{70}$.
\end{enumerate}
\end{multicols}
\vspace{-6pt}

\exo

Écrire les nombres suivants sous la forme $a\sqrt{b}$, où $a$ et $b$ sont des entiers, $b$ étant le plus petit possible.

\vspace{-6pt}
\setlength{\columnsep}{1cm}
\setlength{\columnseprule}{0pt}
\begin{multicols}{2}
\begin{enumerate}
\item $a=\sqrt{50}+4\sqrt{18}-7\sqrt{8}$.
\medskip
\item $b=\sqrt{20}-8\sqrt{45}+2\sqrt{5}$.
\item $c=\sqrt{12}+\sqrt{75}+4\sqrt{300}$.
\medskip
\item $d=5\sqrt{63}-\sqrt{28}+\sqrt{7}$.
\end{enumerate}
\end{multicols}
\vspace{-6pt}

\exo

Écrire les nombres suivants sans radical.

\vspace{-6pt}
\setlength{\columnsep}{1cm}
\setlength{\columnseprule}{0pt}
\begin{multicols}{4}
\begin{enumerate}
\item $a=\sqrt{\dfrac{4}{9}}$.
\item $b=\sqrt{\dfrac{1}{16}}$.
\item $c=\sqrt{\dfrac{49}{25}}$.
\item $d=\dfrac{2}{7}\times\sqrt{\dfrac{49}{64}}$.
\end{enumerate}
\end{multicols}
\vspace{-6pt}

\exo

Écrire les nombres suivants avec un dénominateur entier.

\vspace{-6pt}
\setlength{\columnsep}{1cm}
\setlength{\columnseprule}{0pt}
\begin{multicols}{4}
\begin{enumerate}
\item $a=\dfrac{2}{\sqrt{3}}$.
\item $b=\dfrac{7}{2\sqrt{5}}$.
\item $c=\dfrac{\sqrt{3}}{4\sqrt{2}}$.
\item $d=\dfrac{3\sqrt{3}}{\sqrt{8}}$.
\end{enumerate}
\end{multicols}
\vspace{-6pt}

\exo

Sans les calculer, comparer les nombres suivants. 

\vspace{-6pt}
\setlength{\columnsep}{1cm}
\setlength{\columnseprule}{0pt}
\begin{multicols}{4}
\begin{enumerate}
\item $\sqrt{5}$\; et\; $\sqrt{5,7}\vphantom{\sqrt{\dfrac{5}{2}}}$.
\item $\sqrt{2}$\; et\; $\sqrt{\dfrac{5}{2}}$.
\item $\sqrt{50}$\; et\; $7\vphantom{\sqrt{\dfrac{5}{2}}}$.
\item $\sqrt{\dfrac{10}{3}}$\; et\; $\sqrt{2,7}$.
\end{enumerate}
\end{multicols}
\vspace{-6pt}

\exo

Vrai ou faux ?
\begin{enumerate}
\item La fonction racine carrée est définie sur $\R$.
\item Le nombre $36$ est une solution de l'équation $\sqrt{x}=6$.
\item Le nombre $36$ est la solution de l'équation $\sqrt{x}=6$.
\end{enumerate}

\exo

On a tracé ci-dessous la représentation de la fonction racine carrée sur l'intervalle $\intff{0}{5}$ :
\begin{center}
\def\xmin{-0.5}
\def\xmax{6}
\def\ymin{-0.5}
\def\ymax{3}
\def\xunite{1.0cm}
\def\yunite{1.0cm}
\psset{unit=\xunite, xunit=\xunite, yunit=\yunite, algebraic=true, dimen=middle, dotstyle=*, dotsize=3pt 0, linewidth=0.8pt, arrowsize=3pt 2, arrowinset=0.25, comma=true}
\begin{pspicture*}(\xmin,\ymin)(\xmax,\ymax)
\psgrid[xunit=\xunite, yunit=\yunite, subgriddiv=0, gridlabels=0, gridcolor=lightgray](0,0)(0,0)(\xmax,\ymax)
\psaxes[labelFontSize=\scriptstyle, xAxis=true, yAxis=true, Dx=1, Dy=1, ticksize=-2pt 0, subticks=1]{->}(0,0)(\xmin,\ymin)(\xmax,\ymax)[{\footnotesize $x$},135][{\footnotesize $y$},-45]
\psplot[linewidth=1.2pt, plotpoints=200]{0}{5}{sqrt(x)}
%\psplot[linewidth=1.2pt, plotpoints=200]{\xmin}{\xmax}{2}
\end{pspicture*}
\end{center}
En utilisant ce graphique, résoudre l'inéquation $\sqrt{x}\leqslant 2$.

\exo

Résoudre dans $\R$ les inéquations suivantes.

\vspace{-6pt}
\setlength{\columnsep}{1cm}
\setlength{\columnseprule}{0pt}
\begin{multicols}{4}
\begin{enumerate}
\item $\sqrt{x}<1$.
\medskip
\item $\sqrt{x}>3$.
\item $\sqrt{x} < 9$.
\medskip
\item $\sqrt{x} \geqslant 1$.
\item $2\leqslant \sqrt{x} \leqslant 3$.
\medskip
\item $4< \sqrt{x} \leqslant 5$.
\item $2\sqrt{x}+3 > 5$.
\medskip
\item $-\frac{1}{2}\sqrt{x}-5 < 7$.
\end{enumerate}
\end{multicols}
\vspace{-6pt}

\exo

On considère la proposition : \og Si $x=9$, alors $\sqrt{x}=3$ \fg{}.
\begin{enumerate}
\item Cette proposition est-elle vraie ou fausse ? Justifier.
\item Énoncer la proposition réciproque. Est-elle vraie ou fausse ? Justifier.
\end{enumerate}

\exo

Reprendre l'exercice précédent avec la proposition : \og Si $x=3$, alors $\sqrt{x^2}=3$ \fg{}.

\exo

Un objet tombe en chute libre depuis une hauteur $h$ (en \si{\m}). En supposant qu'il n'est soumis à aucun frottement, sa vitesse (en \si[per-mode=symbol]{\m\per\second}) à l'instant où il touche le sol est donnée par la la relation :
\[v=\sqrt{2gh}, \qouq g=\SI[per-mode=symbol]{9,81}{\m\per\second^2}.\]
\begin{enumerate}
\item Déterminer la vitesse d'impact d'un objet lâché d'une hauteur de \SI{50}{\m}.
\item De quelle hauteur faut-il lâcher l'objet pour que sa vitesse d'impact soit supérieure à \SI[per-mode=symbol]{180}{\km\per\hour} ?
\end{enumerate}

\exo

Un satellite tourne autour de la Terre à l'altitude $h$ (en \si{\km}).

Sa vitesse $v$ (en \si[per-mode=symbol]{\km\per\second}) est alors donnée par la relation :
\[v=\frac{630}{\sqrt{6370+h}}.\]
\begin{enumerate}
\item Quelle est la vitesse d'un satellite situé à \SI[retain-explicit-plus]{35786}{\km} d'altitude ?
\item Quelle est l'altitude d'un satellite qui se déplace à \SI[per-mode=symbol]{2,1}{\km\per\second} ?
\end{enumerate}

